\section{Introduction}

Fine-tuning Large Language Models (LLMs)~\cite{67minaee2024survey, 52stanfordOpportunitiesRisks} is a highly effective way to improve their performance, and add desirable or remove undesirable behaviors~\cite{2ouyang2022training, 13houlsby2019peft, 16howard2018finetuning, 17peters2019finetuning}. Low Rank Adaptation (LoRA)~\cite{1hu2021lora} is one of the most widely adopted methods for fine-tuning LLMs, showing significant promise for enabling smaller, specialized models to outperform larger, more general models on specific tasks, with a fraction of trainable parameters, challenging the notion that bigger general models always outperform smaller ones.

Despite the rapid advancement and release of new base models, such as Gemma~\cite{3gemma}, Llama~\cite{4llama}, and Mistral~\cite{5mistral}, which claim ease of fine-tuning across various tasks, comprehensive evaluations of these models remain scarce. Broad knowledge and reasoning-based benchmarks like MMLU~\cite{28hendrycks2020mmlu} and HellaSwag~\cite{29zellers2019zellers} are commonly used in leaderboards like the Open LLM Leaderboard~\cite{34open-llm-leaderboard}, however, this is not necessarily representative of task-specific performance, before or after fine-tuning. Technical reports~\cite{3gemma, 4llama, 5mistral, 35gpt4, 32gemini} often leave training configurations unspecified, with claims of ease of fine-tuning left unmeasured. While the effectiveness of fine-tuning has been broadly demonstrated~\cite{9kohut2023finetuning, 10zheng2024}, the lack of large-scale experimentation leaves several pivotal questions unanswered, particularly regarding the consistency and predictability of performance improvements through fine-tuning, and the impact of model size, base model, and task complexity.

Evaluations are sensitive to prompting, and there are significant variation in the formulations used in publications and libraries\footnote{https://github.com/openai/simple-evals}. Technical reports often showcase model performance using specialized, dataset-specific prompting strategies such as role-playing prompts (e.g. \textit{"Assume you are an expert"}), \textit{maj@k} voting~\cite{48wangselfconsistency}, varied n-shot~\cite{37song2022evalsurvey}, MedPrompt~\cite{33nori2023generalistvsfinetuned}, and chain-of-thought~\cite{36wei2022cot} prompting. While these methods are intended to highlight the optimal capabilities of models, the use of such diverse prompting techniques can make direct comparisons across models and tasks challenging.

In this work, we seek to bridge these gaps by conducting an extensive analysis of LoRA-based fine-tuning across 10 base models and 31 tasks, totaling 310 LLMs fine-tuned with LoRA. We deliberately maintain that all LLMs are fine-tuned with the same training parameters and emphasize querying with zero or single-shot, completion-style prompts, with simple instructions like \textit{"Solve the following multiple choice problem"}. Altogether, this provides a standardized framework to compare and assess the intrinsic capabilities of different base models when fine-tuned with LoRA under consistent conditions, across specific tasks.

We also aim to explore the viability of serving multiple LoRA models in a real-world production application. LoRAX ~\cite{12predibaseLoRAExchange} enables serving multiple LoRA models simultaneously on a single GPU by leveraging shared base model weights and dynamic adapter loading [12]. We measure latency and concurrency metrics of this library. We use LoRAX to deploy 25 fine-tuned LLM served on a single A100\footnote{https://www.nvidia.com/en-us/data-center/a100/} in the LoRA Land web application. Our successful implementation showcases the economic efficiency of serving multiple LoRA-adapted LLMs for specialized tasks.

Finally, we release all 25 of the fine-tuned models on the LoRA Land web application and their training recipes on (\href{https://huggingface.co/predibase}{Hugging Face}) to allow further analysis and replication by the community.
